% Generated draft: requires mathematical and source review.
\documentclass[11pt]{article}
\usepackage[margin=25mm]{geometry}
\usepackage{fontspec}
\setmainfont{Latin Modern Roman}
\usepackage{amsmath,amssymb,mathtools,unicode-math}
\setmathfont{Latin Modern Math}
\usepackage{enumitem,needspace,xurl,hyperref}
\setlength{\parindent}{0pt}
\setlength{\parskip}{5pt}
\setlength{\emergencystretch}{3em}
\newcommand{\N}{\mathbb{N}}
\newcommand{\Z}{\mathbb{Z}}
\newcommand{\Q}{\mathbb{Q}}
\newcommand{\R}{\mathbb{R}}
\newcommand{\C}{\mathbb{C}}
\newcommand{\F}{\mathbb{F}}
\newcommand{\A}{\mathbb{A}}
\newcommand{\PP}{\mathbb{P}}
\newcommand{\E}{\mathbb{E}}
\newcommand{\eps}{\varepsilon}
\newcommand{\dd}{\,\mathrm{d}}
\begin{document}
\section*{UnsolvedMath ID: 30005395}
\textbf{Source problem code:} catalogue entry
\subsection{Definitions and mathematical statement}
The supplied record formulates a conjectural \(p\)-adic arithmetic Gan--Gross--Prasad height identity. The prime \(p\), the ambient Shimura datum, the subgroup \(H\), and the admissible class of representations \(\pi\) belong to an underlying setup that is not specified in the record. Write \(\operatorname{Sh}(H)\) for the Shimura variety appearing in that setup, \(\pi_f\) for the finite part of \(\pi\), and \(\pi_f^\vee\) for its contragredient. Retain the source notation
\[
Z_\phi:=\int_{\operatorname{Sh}(H)}\phi,
\qquad
Z_{\phi^\vee}:=\int_{\operatorname{Sh}(H)}\phi^\vee.
\]
These are the source's cycle expressions entering a \(p\)-adic height pairing; their construction and target spaces are not supplied, so these formulas are notation rather than complete definitions.

Let \((\ ,\ )_p\) denote the intended \(p\)-adic height pairing. The source states that this pairing depends on an auxiliary splitting of the Hodge filtration and that this choice can be made canonically in the ordinary case. Let \(L_p(s,\pi,R)\) denote the source's \(p\)-adic Rankin--Selberg \(L\)-function, described as interpolating central values of \(L\)-functions twisted by finite order characters. The parameter \(R\), the interpolation normalization, and the derivative convention are not specified. Write \(e_p(\pi)\) for the source's partial Euler factor at \(p\), and \(\alpha_v\) for the local factors occurring in its formula.

\textbf{Requested conjectural identity, conditional on the missing setup and normalizations.} Assume that \(\pi\) is tempered and ordinary at every \(p\)-adic place. For every \(\phi\in\pi_f\) and \(\phi^\vee\in\pi_f^\vee\), the requested formula is
\[
\bigl(Z_\phi,Z_{\phi^\vee}\bigr)_p
=
C\,e_p(\pi)^{-1}
L_p'\!\left(\tfrac12,\pi,R\right)
\prod_v\alpha_v(\phi_v,\phi_v^\vee),
\]
where \(C\) stands for the source's unspecified nonzero explicit normalization factor. It is not a freely chosen factor used to force equality; its value and permitted dependencies must be recovered from the underlying formulation.

The displayed inverse requires \(e_p(\pi)\ne0\), or an explicitly stated replacement when it vanishes; the record supplies neither a nonvanishing assertion nor such a replacement. The local-component notation directly applies to factorizable vectors \(\phi=\bigotimes_v\phi_v\) and \(\phi^\vee=\bigotimes_v\phi_v^\vee\). The source quantifies over all vectors, but does not specify how the product expression extends to arbitrary vectors. The places indexing the product, its normalization, and its meaning are likewise unspecified. Consequently, this preserves the supplied conjectural identity and its quantifiers, but does not constitute a fully determined standalone mathematical conjecture.
\subsection{Short English statement}
The conjecture relates the \(p\)-adic height of Gan--Gross--Prasad cycle expressions to the central derivative of a \(p\)-adic Rankin--Selberg \(L\)-function, multiplied by local factors and an explicit nonzero normalization factor, assuming temperedness and ordinarity at every \(p\)-adic place. The supplied record omits essential geometric data, definitions, and normalization conventions. Its dated curation labels the problem partially solved under additional hypotheses; this draft does not independently verify that assessment.
\subsection{Sources}
\begin{itemize}
\item Local UnsolvedMath registry: supplied complete record, permanent numeric ID \(30005395\), source problem code \texttt{catalogue entry}, title \emph{p-Adic Arithmetic Gan--Gross--Prasad Formula}; collection: Oberwolfach Reports Open Problems. No individual proposer is identified in the supplied record.
\item Supplied report citation: \emph{Arithmetic of Shimura Varieties} (2023), Oberwolfach Reports, pp.~261--326, DOI \texttt{10.4171/owr/2023/5}. \url{https://doi.org/10.4171/owr/2023/5}. The extracted original statement is labelled Conjecture~1; no precise page within the report is supplied.
\item Supplied literature source: Disegni and Zhang, \emph{GGP cycles and p-adic L-functions}. \url{https://arxiv.org/abs/2410.08401}.
\item Supplied literature source: \emph{Oberwolfach Report 5/2023}. \url{https://doi.org/10.4171/owr/2023/5}.
\item Literature assessment dated exactly \texttt{2026-08-22}: the supplied curation states that the \(p\)-adic arithmetic GGP height formula is proved under substantial local, ordinarity, and ramification hypotheses, while the formulation for all tempered ordinary representations remains open. This is supplied dated curation, not an independently verified literature assessment; the precise additional hypotheses are not given in the record.
\item The registry labels the statement \texttt{corrected\_verified} and states that it was reconstructed and checked against the official Oberwolfach report. These are supplied provenance claims; no independent source verification is asserted here.
\end{itemize}
\end{document}
